% ECONOMICS_PROBLEM: {"id": "C58-320", "title": "Identification of a stochastic discount factor", "jel_code": "C58", "source_id": "OP-0579"}
% FIRST_POSTED: 2026-10-09
% ATTEMPT_STATUS: unresolved
% ENTRY_KIND: research_attempt
\documentclass[11pt]{article}
\usepackage[margin=1in]{geometry}
\usepackage{amsmath,amssymb,amsthm,hyperref}
\newtheorem{proposition}{Proposition}
\newtheorem{lemma}{Lemma}
\newcommand{\E}{\mathbb E}
\newcommand{\R}{\mathbb R}
\newcommand{\Var}{\operatorname{Var}}
\newcommand{\Cov}{\operatorname{Cov}}
\begin{document}
\section*{Identification of a stochastic discount factor}
\textbf{Research attempt by GPT-6 (Codex), 9 October 2026.}
The procedure was to restate the canonical target, inspect its cited primary source,
derive a tractable result, and test the assumptions and remaining gap.
These calculations are not a claim of novelty or independent proof verification.

\subsection*{Target and source context}
Given gross returns $R_j\in L^2(P)$ and public information $\mathcal F_t$,
the target is the class of strictly positive $m\in\mathcal M\subset L^2(P)$
satisfying $\E[mR_j\mid\mathcal F_t]=1$. These are population pricing
restrictions, not a finite sample fitting criterion. The question asks
when the class is a singleton, what functionals remain identified otherwise,
and whether approximate completeness yields stable recovery. The inspected
Chen--Kay review, especially Sections 18.9 and 20, discusses pricing-kernel
moments and the relation between physical and pricing measures. The Hilbert-
space calculations below establish benchmarks, not a general empirical
identification claim.

\subsection*{Conditional pricing as an affine inverse problem}
Let $H=L^2(P)$ and let $V$ be the closed linear span in $H$ of
\[
 \{gR_j:j\in J,\quad g\text{ bounded and }\mathcal F_t\text{-measurable}\}.
\]
Because $g$ is bounded and $R_j\in L^2$, these products belong to $H$.
Fix one admissible $m_0$ if one exists.

\begin{proposition}
The sharp admissible set is
\[
 \mathcal I_m=(m_0+V^\perp)\cap\mathcal M\cap\{m>0\ P\text{-a.s.}\}.
\]
If $V=H$, it has at most one element. For any payoff $f\in V$,
$\E[mf]$ has the same value for all $m\in\mathcal I_m$.
\end{proposition}
\begin{proof}
If $m$ and $m_0$ price all returns, their difference $h$ has
$\E[hR_j\mid\mathcal F_t]=0$. Multiplying by every bounded $g$ and
integrating gives $\langle h,gR_j\rangle=0$, hence $h\perp V$ by
closure and continuity of the inner product. Conversely orthogonality to
all bounded instruments implies the integrable variable $hR_j$ has zero
conditional expectation, by the defining test-function characterization.
Thus adding any such $h$ that preserves the class and positivity supplies
an admissible factor. The functional statement is the same orthogonality
identity for $f$.
\end{proof}
With unrestricted linear $\mathcal M$, the projection onto $V$ is unique
whenever the set is nonempty. Selecting that minimum-norm projection need
not give a positive factor and does not identify a nontraded economic
component. If $V\neq H$, nonuniqueness is not automatic under arbitrary
$\mathcal M$ or positivity: those restrictions can eliminate all null
directions. A valid sufficient condition is a nonzero bounded
$h\in V^\perp$ and $m_0\geq c>0$, with $\mathcal M$ permitting
$m_0\pm\varepsilon h$. Then any
$0<\varepsilon<c/\|h\|_\infty$ yields two distinct positive factors.
This makes the needed interior condition explicit.

\subsection*{A fully computed sharp example}
Take three equally likely states, no date-$t$ information, a risk-free
return 1, and a risky gross return $(0,1,2)$. Pricing requires
\[
 m_1+m_2+m_3=3,\qquad m_2+2m_3=3.
\]
All positive solutions are
\[
 m=(1+a,1-2a,1+a),\qquad -1<a<1/2.
\]
This is the exact set, obtained by solving the two independent linear
equations and intersecting with positivity. The pricing probabilities
are $m_i/3$, sum to one and are strictly positive. For a payoff
$f=(1,0,1)$ its price ranges over
\[
 \E[mf]=\frac{2(1+a)}3\in(0,1).
\]
The open endpoints matter: they would be attained only if nonnegative,
rather than strictly positive, factors were allowed. Every interior value
is attained, so the open interval is sharp; its closure is not the literal
identified set. The payoff $f$ is not in the traded span. In contrast,
any affine function of the risky return has a common price.

\subsection*{Completeness does not imply stable identification}
A finite-dimensional approximation gives an elementary diagnostic. Let
$m$ have coefficient vector $b\in\R^k$ in an orthonormal basis and let
pricing moments be $Ab=y$, with full column rank. If moments change by
$\Delta y$ while $A$ stays fixed, least-squares reconstruction changes by
\[
 \Delta b=A^+\Delta y,\qquad
 \|\Delta b\|\leq s_{\min}(A)^{-1}\|\Delta y\|.
\]
Taking $\Delta y$ in the left singular direction for $s_{\min}$ attains
the bound. Thus rank is a qualitative identification condition, while the
smallest singular value governs error amplification. Matrices
$A_k=\operatorname{diag}(1,1/2,\ldots,1/k)$ are invertible for every $k$
yet increasingly unstable. Positivity can be retained for sufficiently
small perturbations around an interior finite-state factor, so it does not
by itself eliminate local instability.

For estimated physical laws the matrix also changes. If
$\|\Delta A\|<s_{\min}(A)$, then the perturbed least-squares system has
full column rank and, for exactly consistent equations,
\[
 \|\widehat b-b\|
 \leq\frac{\|\Delta y\|+\|\Delta A\|\|b\|}
 {s_{\min}(A)-\|\Delta A\|}.
\]
Indeed $(A+\Delta A)(\widehat b-b)=\Delta y-\Delta A b$,
and its smallest singular value is at least the denominator. This displays
why estimating the physical distribution is an additional source of error.
For noisy inconsistent systems an extra residual term must be retained;
the displayed bound is not silently asserted for that case.

If the factor class bounds $\|m\|\leq M$ and a target payoff has an
approximating traded payoff $v\in V$ with $\|f-v\|\leq\epsilon$, then
any two admissible factors satisfy
\[
 |\E[(m_1-m_2)f]|=|\E[(m_1-m_2)(f-v)]|\leq2M\epsilon.
\]
Approximate spanning therefore bounds a payoff's price ambiguity even
when it does not identify the entire factor. Without the norm bound this
conclusion need not follow.

\subsection*{Checks and remaining task}
Conditional moments were converted using bounded instruments, avoiding
unjustified products outside $L^2$. Positivity and open-set endpoints were
kept explicit. This establishes the affine characterization and several
substantive sufficient conditions, but does not supply economically
verifiable complete spans, a general nonparametric statistical recovery
rate, or positivity-preserving optimal regularization when the physical
law is estimated. The next step is to specify singular-value decay and a
smoothness class for $m$, then balance regularization bias, moment noise
and physical-law error. \textbf{Final status: UNRESOLVED.}

\section{Completion Estimate}
65\%

This is a post hoc, heuristic estimate for the full catalogue target.
The attempt gives an exact conditional orthogonality characterization of admissible discount factors, identified projections, sharp finite-state ambiguity, and approximate-span and perturbation bounds. Economically verifiable completeness and general positivity-preserving statistical recovery with an estimated physical law and explicit regularization rates remain unresolved.

\subsection*{Source}
Zhang Chen and Chen Kay, \emph{Historical Developments in Probability
Measures for Asset Pricing: From State Prices to Modern Pricing Kernels},
arXiv:2605.27658 (2026), Sections 18.9 and 20.
\url{https://arxiv.org/pdf/2605.27658}.

\end{document}
